\documentclass[letter,10pt]{article}
\usepackage[utf8]{inputenc}

\usepackage{TLCresume}

\def\name{Jhulen Mallo}
\def\phone{(+591) 79577091}
\def\city{La Paz, Bolivia}
\def\email{jmalloq@gmail.com}
\def\LinkedIn{jhulen-mallo}
\def\github{}
\def\role{Ingeniero de Datos | Certif{}icado en Google Cloud}
\hypersetup{pdftitle={Jhulen Mallo - Ingeniero de Datos | Certif{}icado en Google Cloud},pdfauthor={Jhulen Mallo},pdfsubject={Experiencia y certificaciones profesionales}}

\RequirePackage{fancyhdr}
\fancypagestyle{fancy}{%
\fancyhf{}
\lhead{\phone \\
        \city \\
	    \href{mailto:\email}{\email}}
	\chead{%
	    \centering {\Huge \skills \name \vspace{.25em}} \\
	    {\color{highlight} \Large{\role}}}%
	    \rhead{\href{https://jhulenmallo.github.io/cv/}{Sitio web}\\
            \href{https://jhulenmallo.github.io/cv/}{jhulenmallo.github.io/cv/} \\
	    \href{https://www.linkedin.com/in/\LinkedIn}{linkedin.com/in/\LinkedIn}}
\renewcommand{\headrulewidth}{1pt}
\renewcommand{\headrule}{\hbox to\headwidth{%
  \color{highlight}\leaders\hrule height \headrulewidth\hfill}}
}
\pagestyle{fancy}

\setlength{\headheight}{58pt}
\setlength{\headsep}{5pt}
\begin{document}
Ingeniero de Datos con experiencia en banca, integración de información y automatización de procesos con Python, SQL y Apache Air{f}low. He liderado la orquestación de flujos desde un Data Lake hacia sistemas de negocio y el desarrollo de APIs sobre OpenShift/Kubernetes. Combino ingeniería y analítica con formación en informática y estadística para construir soluciones mantenibles que apoyen la toma de decisiones.
\vspace{0.3em}

\section{Habilidades}
\begin{tabular}{p{9.5em} p{1em} p{\dimexpr\textwidth-10.5em\relax}}
\skills{Ingeniería de Datos} & & Python, SQL, ETL, Apache Air{f}low, DAGs, Data Lake, SQL Server, PostgreSQL. \\
\skills{Cloud y Despliegue} & & Google Cloud Platform, Azure App Services, Azure Blob Storage, Docker, Kubernetes, OpenShift, Git. \\
\skills{Servicios y BI} & & FastAPI, GraphQL, Clean Architecture, Arquitectura Hexagonal, Power BI, Power Automate. \\
\skills{IA y Formación} & & RAG, LangChain, ChromaDB, Gemini; formación en Spark, machine learning y MLOps. \\
\end{tabular}


\section{Experiencia Profesional}
\subsection{{DATA ENGINEER \hfill Julio 2026 --- Presente}}
\subtext{Banco Bisa S.A. \hfill La Paz, Bolivia}
\begin{zitemize}
\item Mantenimiento y optimización de pipelines y procesos ETL mediante Apache Air{f}low y Docker.
\item Desarrollo y mantenimiento de microservicios y APIs en Python con FastAPI y GraphQL, aplicando Clean Architecture y Arquitectura Hexagonal.
\item Participación en el diseño de un Portal Analítico Empresarial, con foco en gobernanza, reproducibilidad y auditabilidad de la información.
\item Desarrollo de un prototipo RAG sobre documentación normativa interna con LangChain, ChromaDB y Gemini.
\end{zitemize}

\subsection{{DATA ANALYST \hfill Junio 2025 --- Julio 2026}}
\subtext{Banco Bisa S.A. \hfill La Paz, Bolivia}
\begin{zitemize}
\item Lideré la integración de datos del Data Lake con el nuevo CRM mediante procesos batch y de baja latencia.
\item Lideré la primera implementación de Apache Air{f}low del área: arquitectura de orquestación, plantillas de DAGs y estándares.
\item Desarrollé la primera capa de APIs en Python del área con FastAPI y GraphQL, desplegada sobre OpenShift/Kubernetes.
\item Colaboré con Arquitectura de Datos y Sistemas en el diseño de iniciativas de modernización de datos.
\end{zitemize}

\subsection{{RISK DATA ANALYST \hfill Enero 2024 --- Mayo 2025}}
\subtext{Banco Bisa S.A. \hfill La Paz, Bolivia}
\begin{zitemize}
\item Extracción, transformación y preparación de datos con Python y SQL Server para modelos y dashboards de Riesgos; automatización de reportes con Power BI y Power Automate.
\item Análisis de cartera y comportamiento de clientes: scoring, Probabilidad de Default, mora, cosechas, churn y Customer Lifetime Value.
\end{zitemize}

\subsection{{FULL STACK DEVELOPER \hfill Noviembre 2023 --- Diciembre 2025}}
\subtext{Santillana Bolivia (Consultoría remota) \hfill La Paz, Bolivia}
\begin{zitemize}
\item Desarrollo de APIs REST e interfaces web para módulos educativos con .NET, SQL Server, JavaScript y Bootstrap, desplegadas sobre Azure App Services y Blob Storage.
\end{zitemize}


\section{Educación}

\skills{Licenciatura en Informática (Mención Ciencias de la Computación)}, \textit{Universidad Mayor de San Andrés}	\hfill 2018 --- 2024

\skills{Licenciatura en Estadística (enfoque en Data Science, en curso)}, \textit{Universidad Mayor de San Andrés} \hfill	2023 --- Presente


\section{Actividades}
Ayudante de Cátedra (UMSA): Programación, POO, Análisis Numérico y Laboratorio de Computación. \hfill 2020 --- 2023

Premios: Primer Lugar Feria OPEN-INFO (2021) y Certif{}icado de Rendimiento Académico (2021).


\section{Certif{}icaciones y Formación Seleccionada}
\begin{zitemize}
    \item \href{https://www.credly.com/badges/adc543d1-0d0a-42cf-9b76-f09412427cd1}{Associate Cloud Engineer} - \textit{Google Cloud} \hfill Feb 2026
    \item \href{https://www.coursera.org/account/accomplishments/specialization/94N7VQDPURXZ}{Data Engineering, Big Data, and Machine Learning on GCP} - \textit{Google Cloud - Coursera} \hfill Sep 2023
    \item \href{https://www.coursera.org/account/accomplishments/specialization/TQAEV5VXPVRC}{MLOps | Machine Learning Operations} - \textit{Duke University - Coursera} \hfill Sep 2023
    \item \href{https://www.credly.com/badges/6cf3e9ab-4273-437c-b1a1-de7fd75c58f8}{Generative AI Leader Certification (Spanish)} - \textit{Google Cloud} \hfill Nov 2026
    \item \href{https://www.coursera.org/account/accomplishments/specialization/YV2XYVC2GGS3}{Mathematics for Machine Learning and Data Science} - \textit{DeepLearning.AI - Coursera} \hfill Oct 2023
    \item \href{https://www.coursera.org/account/accomplishments/verify/9RVETDFCQ98N}{Generative AI with Large Language Models} - \textit{DeepLearning.AI \& Amazon Web Services - Coursera} \hfill Sep 2023
    \item \href{https://www.coursera.org/account/accomplishments/specialization/4NN6BWAJ6Y9C}{Deep Learning} - \textit{DeepLearning.AI - Coursera} \hfill Oct 2023
\end{zitemize}

\end{document}